% This file is generated from the corresponding Obsidian note.
% Source: Teoricos/GS - Teo6.md
\chapter{El espacio tangente y la diferencial}
\label{ch:teo-6}
\section{El espacio tangente de una carta}
\begin{bookremark}{El espacio tangente de una carta}
Si $(U,\varphi=(x_1,\dots,x_n))$ es una carta suave de $p$ en $M$, entonces también tiene sentido considerar el espacio tangente $T_pU$, ya que $U$ es por sí mismo una variedad suave.

En $T_pU$ aparece automáticamente la base

\[
\left\{\left.\frac{\partial}{\partial x_i}\right|_p\right\}_{i=1}^n.
\]
No se puede decir literalmente que $T_pM$ y $T_pU$ sean el mismo conjunto, porque los elementos de $T_pM$ actúan sobre funciones en $C^\infty(M)$ y los de $T_pU$ actúan sobre funciones en $C^\infty(U)$.

Sin embargo, son isomorfos de forma natural: al restringir funciones globales a $U$ y al extender funciones locales cerca de $p$, los mismos vectores coordenados describen ambos espacios.

Una forma elegante de evitar este abuso es definir vector tangente como una derivación sobre el álgebra de gérmenes $C_p^\infty(M)$.

\end{bookremark}
\section{Diferencial de una función suave}
\begin{bookdefinition}{Diferencial de una función suave}
Sean $M$ y $N$ variedades suaves y $F:M\to N$ una función suave. Para cada $p\in M$ se define la función

\[
dF_p:T_pM\to T_{F(p)}N,
\]
llamada \textbf{diferencial de $F$ en $p$}, por

\[
\bigl(dF_p(v)\bigr)(g)=v(g\circ F),\qquad g\in C^\infty(N).
\]
Está bien definida porque si $g\in C^\infty(N)$, entonces $g\circ F\in C^\infty(M)$.

\end{bookdefinition}
\begin{bookexercise}{Observación sobre $dF_p$}
$dF_p$ es lineal y el vector tangente $dF_p(v)$ es lineal y vuelve a satisfacer la regla de Leibniz.

\end{bookexercise}
\section{Regla de la cadena}
\begin{bookproposition}{Regla de la cadena}
Si $F:M\to N$ y $G:N\to P$ son suaves, entonces

\[
d(G\circ F)_p=(dG)_{F(p)}\circ dF_p.
\]
\end{bookproposition}
\begin{bookproof}{}
\begin{enumerate}
\item Sea $v\in T_pM$ y $h\in C^\infty(P)$. Entonces como la composición es asociativa
\end{enumerate}
\[
\bigl(d(G\circ F)_p(v)\bigr)(h)=v(h\circ (G\circ F))=v(h\circ G\circ F)
\]
\begin{enumerate}[start=2]
\item Por otro lado, mirando $dF_{p}(v)$ como un vector tangente y aplicandolo a $h\circ G$
\end{enumerate}
\[
\bigl(((dG)_{F(p)}\circ dF_p)(v)\bigr)(h)=((dG)_{F(p)}(dF_{p}(v)))(h)=(dF_p(v))(h\circ G)=v((h\circ G)\circ F)=v(h\circ G\circ F).
\]
\begin{enumerate}[start=3]
\item Como ambas derivaciones actúan igual sobre toda función $h\in C^\infty(P)$, son iguales.
\item Y luego como ambas actuan sobre todo $v\in T_{p}M$ son iguales
\end{enumerate}
\bookqed
\end{bookproof}
\section{Diferencial de un difeomorfismo}
\begin{bookcorollary}{El diferencial de un difeomorfismo es un isomorfismo}
Si $F:M\to N$ es un difeomorfismo, entonces

\[
(dF)_p:T_pM\to T_{F(p)}N
\]
es un isomorfismo, y su inversa es

\[
\bigl((dF)_p\bigr)^{-1}=(dF^{-1})_{F(p)}.
\]
\end{bookcorollary}
\begin{bookproof}{}
Como $\operatorname{Id}_M=F^{-1}\circ F$, por la regla de la cadena se obtiene

\[
\operatorname{Id}_{T_pM}=d(\operatorname{Id}_M)_p=(dF^{-1})_{F(p)}\circ(dF)_p.
\]
Análogamente, usando $\operatorname{Id}_N=F\circ F^{-1}$,

\[
\operatorname{Id}_{T_{F(p)}N}=d(\operatorname{Id}_N)_{F(p)}=(dF)_p\circ(dF^{-1})_{F(p)}.
\]
Luego $(dF)_p$ es invertible y su inversa es $(dF^{-1})_{F(p)}$.

\bookqed
\end{bookproof}
\phantomsection\label{blk:3d3615}
\section{Identificación entre $T_pU$ y $T_pM$}
\begin{bookproposition}{La inclusión induce un isomorfismo}
Sea $U\subseteq M$ un abierto, $p\in U$ e $i:U\hookrightarrow M$ la inclusión (funcion suave). Entonces

\[
(di)_p:T_pU\to T_pM
\]
es un isomorfismo.

\end{bookproposition}
\begin{bookproof}{}
\begin{enumerate}
\item Como $\dim T_pU=\dim T_pM=\dim M$, basta ver que $(di)_p$ es inyectiva.
\item Sea $v\in T_pU$ tal que $(di)_p(v)=0$. Queremos probar que $v=0$.
\item Tomemos $f\in C^\infty(U)$ arbitraria. Por el lema de extensión, existe una extensión suave $\widetilde f\in C^\infty(M)$ definida en un abierto de $p$ y tal que $\widetilde f=f$ cerca de $p$.
\item Entonces
\end{enumerate}
\[
0=((di)_p(v))( \widetilde f )=v(\widetilde f\circ i)=v(f)
\]
porque $\widetilde f\circ i=f$ en un entorno de $p$. 5. Como esto vale para toda $f\in C^\infty(U)$, se concluye que $v=0$. Por lo tanto, $(di)_p$ es inyectiva y en consecuencia un isomorfismo.

\bookqed
\end{bookproof}
\begin{bookremark}{Vectores coordenados y la inclusión}
Si $(U,\varphi=(x_1,\dots,x_n))$ es una carta suave alrededor de $p$, el isomorfism

\[
(di)_p:T_pU\to T_pM
\]
\end{bookremark}
envía cada vector coordenado de $T_pU$ en el vector coordenado correspondiente de $T_pM$

\[
(di)_p\left(\left.\frac{\partial}{\partial x_i}\right|_p\right)=\left.\frac{\partial}{\partial x_i}\right|_p
\]
. Esto permite identificar sin ambigüedad $T_pU$ con $T_pM$.

\begin{bookremark}{Comentario de Lee}
Lee señala que esta identificación es canónica e independiente de elecciones: un vector tangente en $p$ actúa sólo sobre el germen local de las funciones cerca de $p$. Por eso no hay diferencia esencial entre pensarlo como derivación sobre funciones definidas en $U$ o sobre funciones definidas en todo $M$.

\end{bookremark}
\section{Ejercicio}
\begin{bookexercise}{Diferencial nulo y funciones constantes}
Sean $M$ y $N$ variedades suaves y sea $F:M\to N$ una función suave, con $M$ conexa. Probar que

\[
(dF)_p:T_pM\to T_{F(p)}N
\]
es la transformación nula para todo $p\in M$ si y sólo si $F$ es constante.

\end{bookexercise}
\section{Cambio de coordenadas}
\begin{bookproposition}{Matriz de cambio de base entre cartas}
Sean

\[
(U,\varphi=(x_1,\dots,x_n)),\qquad (V,\psi=(y_1,\dots,y_n))
\]
dos cartas suaves alrededor de $p$.

Las bases coordenadas

\[
\mathcal B_1=\left\{\left.\frac{\partial}{\partial x_1}\right|_p,\dots,\left.\frac{\partial}{\partial x_n}\right|_p\right\},\qquad\mathcal B_2=\left\{\left.\frac{\partial}{\partial y_1}\right|_p,\dots,\left.\frac{\partial}{\partial y_n}\right|_p\right\}
\]
Entonces la matriz $C$ cambio de base es

\[
\bigl[C\bigr]_{\mathcal B_2}^{\mathcal B_1}=J(\varphi\circ\psi^{-1})(\psi(p))
\]
\end{bookproposition}
\begin{bookproof}{}
\begin{enumerate}
\item Escribimos
\end{enumerate}
\[
\left.\frac{\partial}{\partial y_j}\right|_p=\sum_{i=1}^n a_{ij}\left.\frac{\partial}{\partial x_i}\right|_p
\]
\begin{enumerate}[start=2]
\item Para hallar los coeficientes, aplicamos ambos lados a la función coordenada $x_i$:
\end{enumerate}
\[
a_{ij}=\left.\frac{\partial}{\partial y_j}\right|_p(x_i)
\]
\begin{enumerate}[start=3]
\item Por definición de vector coordenado,
\end{enumerate}
\[
a_{ij}=\left.\frac{\partial}{\partial t_j}\right|_{\psi(p)}(x_i\circ\psi^{-1})
\]
donde $(t_1,\dots,t_n)$ son las coordenadas usuales en $\mathbb{R}^n$. 4. Pero $x_i=\pi_i\circ\varphi$, donde $\pi_i:\mathbb{R}^n\to\mathbb{R}$ es la proyección sobre la $i$-ésima coordenada. Luego

\[
x_i\circ\psi^{-1}=\pi_i\circ\varphi\circ\psi^{-1}
\]
\begin{enumerate}[start=5]
\item Por lo tanto
\end{enumerate}
\[
a_{ij}=\left.\frac{\partial}{\partial t_j}\right|_{\psi(p)}\bigl(\pi_i\circ\varphi\circ\psi^{-1}\bigr)
\]
que es exactamente la entrada $(i,j)$ de la jacobiana de $\varphi\circ\psi^{-1}$ en $\psi(p)$. 6. Así, la matriz de cambio de base es

\[
C(B_{2},B_{1})=J(\varphi\circ\psi^{-1})(\psi(p))
\]
\bookqed
\end{bookproof}
